\chapter{Dual-Core Bring-up Test}
\label{chap:fpgasteel-dualcore-bringup}

\section{About this project}
\label{sec:fpgasteel-dualcore-about}

\textit{``Dual-Core Bring-up Test''} (under \repopath{sw/dualcore_bringup_test}) is the smallest possible test that gets both \term{Cortex-A9} cores on the \gterm{cyclonevsoc}{Cyclone V SoC}\footnote{\glsentrydesc{cyclonevsoc}} running \gterm{smp}{at the same time}: \term{CPU0} turns on an alternating \term{LED} pattern and then boots \term{CPU1}; \term{CPU1}, once running, shows 123 on the \term{7-segment display} and loops forever. There is no camera, no \gterm{vga}{VGA}, no interrupts, no \gterm{cpp}{C++}: the whole project is two small \term{C} functions and one assembly stub, on purpose, so that the only thing being exercised is the act of bringing \term{CPU1} up at all. \nameref{chap:fpgasteel-dualcore-forwarding} builds on top of this same bring-up sequence to add cache, \gterm{scu}{SCU}, \gterm{mmu}{MMU} configuration, and actual work dispatched to \term{CPU1}; this project is its minimal foundation.

The sequence used to bring \term{CPU1} up is not documented step-by-step anywhere in Intel's own \term{TRM} \cite{CV-TRM} or \gterm{hwlib}{hwlib}: the \term{TRM} explains the reset architecture in general terms, but not as a ready-made recipe. The exact order of operations implemented here, reset-assert, trampoline write, start-address write, cache clean, reset-deassert, is the same sequence the \term{Linux} kernel uses to bring up its second core, implemented in \extpath{arch/arm/mach-socfpga/platsmp.c}. This repository does not vendor a copy of the \term{Linux} kernel source; the sequence in \repopath{main.c} was arrived at by using \gterm{claudecode}{Claude Code} (Anthropic's \term{AI} coding assistant) to read and explain the relevant parts of that kernel file.

\section{Prerequisites}
\label{sec:fpgasteel-dualcore-prerequisites}

Same toolchain, \term{SD card} and \term{Eclipse}/\gterm{openocd}{OpenOCD} debug setup as \nameref{chap:fpgasteel-hello-world}, including the \repopath{boot.script.jtag_debug} script, the same one used by every example in this project.

\section{FPGA Hardware Involved}
\label{sec:fpgasteel-dualcore-hardware}

Only two \term{Qsys}/\gterm{platformdesigner}{Platform Designer} peripherals are used, both reached through the lightweight \gterm{hps}{HPS}-to-\gterm{fpga}{FPGA} bridge:

\begin{itemize}
    \item \hw{hex_display_controller}, the same core and address documented in Paragraph~\ref{subsec:fpgasteel-hex-display}. \term{CPU1} writes 123 to it once it starts running.
    \item An \term{LED} \gterm{pio}{PIO}\footnote{\glsentrydesc{pio}} core, driving the \term{DE1-SoC}'s ten \term{LED}s: \term{CPU0} writes 0x2AA (binary 1010101010) to it once, right after boot, to show it is alive.
\end{itemize}

Neither the camera capture nor the \gterm{vga}{VGA} output datapaths are used here.

\section{Operating principles}
\label{sec:fpgasteel-dualcore-operating-principles}

\subsection{CPU1 starts held in reset}
\label{subsec:fpgasteel-dualcore-cpu1-reset}

On the \term{Cyclone V HPS}, only one core ever comes out of reset by itself. This behavior is documented in \cite[Chapter 10, ``\term{Cortex-A9} Microprocessor Unit Subsystem'', p.~10-6]{CV-TRM}, which states:

\begin{quote}
    ``When a cold or warm reset is issued in the \gterm{mpu}{MPU}\footnote{\glsentrydesc{mpu}}, \term{CPU0} is released from reset automatically. If \term{CPU1} is present, then its reset signals are left asserted when a cold or warm reset is issued. After \term{CPU0} comes out of reset, it can deassert \term{CPU1}'s reset signals by clearing the \term{CPU1} bit in the \gterm{mpu}{MPU} Module Reset (\hw{mpumodrst}) register in the \gterm{resetmanager}{Reset Manager}\footnote{\glsentrydesc{resetmanager}}.''
\end{quote}

The \gterm{resetmanager}{Reset Manager} chapter \cite[Chapter 4, ``Reset Sequencing'', p.~4-11]{CV-TRM} states the same fact from the reset controller's own point of view:

\begin{quote}
    ``After the reset manager releases the \gterm{mpu}{MPU} subsystem from reset, \term{CPU1} is left in reset and \term{CPU0} begins executing code from the reset vector address. Software is responsible for deasserting \term{CPU1}[...]''.
\end{quote}

So \term{CPU1} is not merely idle, parked, or spinning after boot: it is held in actual hardware reset, executing nothing, from the moment the board powers on until some piece of software running on \term{CPU0}, \gterm{uboot}{U-Boot}, \term{Linux}, or, as here, this bare-metal test, explicitly clears the \term{CPU1} bit of the \hw{mpumodrst} register (\hw{RSTMGR_MPUMODRST}, 0xFFD05010, \repopath{main.c}, Paragraph~\ref{subsec:fpgasteel-dualcore-main-c}). Nothing in \gterm{uboot}{U-Boot}'s own boot process does this on its own; \gterm{uboot}{U-Boot} itself runs single-core and simply never touches that bit, which is exactly why \term{CPU1} is still sitting in reset, untouched, by the time this project's \ccode{main()} gets control.

\subsection{The boot address and the remap trick}
\label{subsec:fpgasteel-dualcore-boot-address}

\ccode{main()}, running on core 0, asserts \term{CPU1}'s reset again first (\hw{RSTMGR_MPUMODRST} = 0x2), even though it should already be asserted at this point; this makes the bring-up idempotent regardless of whatever state \term{CPU1} was actually left in, rather than assuming the hardware default was never touched.

It then writes a two-word trampoline directly to physical address 0x0:

\begin{lstlisting}[language=C]
trampoline[0] = 0xE51FF004; /* ldr pc, [pc, #-4] */
trampoline[1] = (uint32_t)&_start_core1; /* target address */
\end{lstlisting}

0xE51FF004 is the \term{ARM} encoding of {\catcode`\#=12 \ccode{ldr pc, [pc, #-4]}}\footnote{\texttt{ldr} = load, \texttt{pc} = program counter, \texttt{\#-4} = next instruction.}: loaded at address 0, this instruction reads the 32-bit word stored 4 bytes behind the current program counter, i.e.\ the second trampoline word, and branches to it. This two-instruction idiom is the standard way to perform a long, absolute jump to an address that a plain relative branch instruction cannot reach on its own.

This only works because address 0x0 is not, on this board at this moment, what it is by default. The system interconnect address map \cite[Chapter 8, Figure 8-4, p.~8-6]{CV-TRM} shows that for the \gterm{mpu}{MPU}, address 0x0 normally maps to the on-chip \gterm{bootrom}{Boot ROM}, not \term{SDRAM}. \gterm{uboot}{U-Boot}, however, reconfigures the system's interconnect remap register early in its own boot process so that address 0x0 maps to \term{SDRAM} instead (needed for \gterm{uboot}{U-Boot}'s own exception vectors); this project's trampoline write relies entirely on that remap already being in place by the time \ccode{main()} runs.

\term{CPU1}'s start address register, \hw{cpu1startaddr}, is written too, with the same target address. According to \cite[Chapter 6, ``\term{L3 Interconnect}'', p.~6-7]{CV-TRM}, this register is what the \gterm{bootrom}{Boot ROM}'s own reset handler reads to know where to jump, in the ordinary case where \term{CPU1} actually restarts execution from the \gterm{bootrom}{Boot ROM}. Since the remap above means \term{CPU1} will execute the trampoline at 0x0 directly instead, this write is redundant in this specific setup.

\subsection{Cache maintenance for the trampoline}
\label{subsec:fpgasteel-dualcore-cache-maintenance}

Before releasing \term{CPU1}, \ccode{main()} explicitly cleans the two trampoline words from the \term{CPU}'s own caches back to \term{DDR}: the \term{L1} data cache via the \ccode{DCCMVAC} coprocessor operation (\ccode{mcr p15, 0, Rd, c7, c10, 1}, once per word), then the shared \term{PL310} \term{L2} controller via its ``Clean Line by PA'' and ``Cache Sync'' registers (0xFFFEF7B0, 0xFFFEF730). Both are followed by a \ccode{dsb} to guarantee the writes have actually completed before reset is deasserted.

This project's own boot script (\repopath{boot.script.jtag_debug}, Section~\ref{sec:fpgasteel-dualcore-prerequisites}) leaves the caches disabled by the time \ccode{main()} runs, which would make this cleaning step a no-op here: cleaning an address that was never cached has nothing to write back. The sequence is included anyway because it is copied unconditionally from \term{Linux}'s \texttt{platsmp.c}, where it runs with the primary core's caches fully enabled and is not optional; keeping the same steps here, rather than trimming what looks redundant in this particular context, is what makes the sequence a faithful, directly-comparable copy of the mechanism \term{Linux} itself relies on.

\subsection{Releasing CPU1}
\label{subsec:fpgasteel-dualcore-releasing-cpu1}

The last step is simply clearing the reset bit again (\hw{RSTMGR_MPUMODRST} = 0x0). \term{CPU1} immediately fetches from address 0x0, runs the two-word trampoline from Paragraph~\ref{subsec:fpgasteel-dualcore-boot-address}, and lands in \ccode{_start_core1}.

\section{Code structure}
\label{sec:fpgasteel-dualcore-code-structure}

\subsection{main.c}
\label{subsec:fpgasteel-dualcore-main-c}

Plain \term{C}, no classes. Source: \repopath{main.c}.

\begin{itemize}
    \item \ccode{main()} (\term{CPU0} entry point): turns on the \term{LED} pattern, shows dashes on the hex display, then runs the five-step \term{CPU1} bring-up sequence from Section~\ref{sec:fpgasteel-dualcore-operating-principles} (assert reset, write trampoline, write \hw{cpu1startaddr}, clean caches, deassert reset), waits, and loops forever.
    \item \ccode{main1()} (\term{CPU1} entry point, called from \repopath{startup_core1.S}): writes 123 to the hex display and loops forever. This is the only code that ever runs on \term{CPU1} in this project.
\end{itemize}

\subsection{startup\_core1}
\label{subsec:fpgasteel-dualcore-startup-core1}

\term{CPU1}'s own minimal entry stub, since \term{CPU1} starts with no stack and no runtime set up at all, unlike \term{CPU0} (which inherits \gterm{uboot}{U-Boot}'s stack until its own \term{C} runtime takes over). Source: \repopath{startup_core1.S}.

\begin{itemize}
    \item \ccode{_start_core1}: sets the stack pointer to \ccode{_stack1}, then calls \ccode{main1()}. If \ccode{main1()} ever returned, which it does not, execution would fall into a \ccode{wfe}/branch-to-self trap loop immediately after.
\end{itemize}

\subsection{Linker script (ldscripts/cvav-ddr.ld)}
\label{subsec:fpgasteel-dualcore-linker-script}

Identical to the single-core project template (the same one \nameref{chap:fpgasteel-simple-forwarding} uses) except for one change: the template's single stack region is split into two, so that \term{CPU0} and \term{CPU1} never share, or overlap, the memory backing their stacks.

\begin{table}[htbp]
    \centering
    \begin{tabular}{@{}l p{4.5cm} p{5.5cm}@{}}
        \toprule
        & \textbf{Single-core template} & \textbf{This project} \\
        \midrule
        \term{CPU0} stack & \ccode{.stack} at 0x80000\newline{}512 KB\newline{}symbol \ccode{_stack} & \ccode{.stack0} at 0x1000000\newline{}16 MB\newline{}symbol \ccode{_stack0} \\
        \lightmidrule
        \term{CPU1} stack & (none) & \ccode{.stack1} at 0x2000000\newline{}16 MB\newline{}symbol \ccode{_stack1} \\
        \bottomrule
    \end{tabular}
    \caption{Stack layout: single-core template vs.\ this project.}
    \label{tab:fpgasteel-dualcore-stack-layout}
\end{table}

\ccode{_stack1} is not a cosmetic addition: \repopath{startup_core1.S} loads it directly into stack pointer before calling \ccode{main1()}, so this linker change is what actually gives \term{CPU1} a valid, dedicated stack to run on at all. Without it, \term{CPU1} would start with whatever value the stack pointer happened to already contain, most likely 0x0, the address that has just been overwritten with the trampoline.

\section{How it works}
\label{sec:fpgasteel-dualcore-how-it-works}

The full sequence, \term{CPU0}'s and \term{CPU1}'s sides together:

\begin{enumerate}
    \item \gterm{openocd}{OpenOCD} debugger halts \term{CPU0} over \gterm{jtag}{JTAG} and loads the \term{ELF} file. \term{CPU1} is still held in hardware reset and has not executed anything at all yet.
    \item \term{CPU0} runs \ccode{main()}: \term{LED}s turn on, dashes show on the hex display.
    \item \term{CPU0} runs the five-step bring-up sequence (Paragraphs~\ref{subsec:fpgasteel-dualcore-boot-address}--\ref{subsec:fpgasteel-dualcore-releasing-cpu1}).
    \item \term{CPU1} comes out of reset, executes the trampoline at 0x0, and lands in \ccode{_start_core1}, which sets up its stack and calls \ccode{main1()}.
    \item \term{CPU1} writes 123 to the hex display and loops forever; \term{CPU0}, after a short busy-wait, also loops forever.
\end{enumerate}

If the hex display shows 123, both cores are confirmed running independently, which is the entire point of this test.